% Display math whose body is only a quantikz circuit of quantikz 1.0 and
% later: in \[ \], $$ $$, displaymath, equation, equation*, align*,
% gather*, a theorem, and a proof with \qedhere, in a \resizebox and a
% \scalebox, with options, a \label before and after it, a \tag, a period
% or a comma after it, and the common gate commands. Each case follows a
% paragraph that names it. A circuit in a figure and ordinary display
% math stay as they are.
\documentclass{article}
\usepackage{amsmath}
\usepackage{amsthm}
\usepackage{graphicx}
\usepackage{quantikz}
\newtheorem{theorem}{Theorem}
\begin{document}
Bracket quantikz.
\[ \begin{quantikz} \lstick{$|0\rangle$} & \gate{H} & \qw \end{quantikz} \]

Dollars quantikz.
$$ \begin{quantikz} \lstick{$|1\rangle$} & \gate{X} & \qw \end{quantikz} $$

Displaymath quantikz.
\begin{displaymath}
\begin{quantikz} & \gate{Z} & \meter{} \end{quantikz}
\end{displaymath}

Equation quantikz period.
\begin{equation}
\begin{quantikz}[row sep=0.3cm, column sep=0.4cm]
\lstick{$a$} & \ctrl{1} & \qw \\
\lstick{$b$} & \targ{} & \qw
\end{quantikz}.
\label{eq:circuit}
\end{equation}

Equation star quantikz.
\begin{equation*}
\begin{quantikz} & \gate{T} & \gate{S} & \qw \end{quantikz}
\end{equation*}

Align star quantikz.
\begin{align*}
\begin{quantikz}
& \ctrl{1} & \gate{H} \\
& \targ{} & \qw
\end{quantikz}
\end{align*}

Gather star quantikz.
\begin{gather*}
\begin{quantikz} & \swap{1} & \qw \\ & \targX{} & \qw \end{quantikz}
\end{gather*}

Bracket resizebox quantikz.
\[ \resizebox{0.25\textwidth}{!}{%
\begin{quantikz}
\lstick{$|0\rangle$} & \gate{H} & \gate[2]{U_f} & \gate{H} & \meter{} \\
\lstick{$|1\rangle$} & \gate{H} & & \qw & \qw
\end{quantikz}} \]

Bracket scalebox quantikz comma.
\[ \scalebox{0.8}{\begin{quantikz} & \gate{R_x(\theta)} & \qw \end{quantikz}}\, , \]

Bracket wire types quantikz.
\[ \begin{quantikz}[wire types={q,c}]
\lstick{$q$} & \gate{H} \slice{$t_1$} & \meter{} \\
\lstick{$c$} & & \setwiretype{c}
\end{quantikz} \]

Bracket bundle quantikz.
\[ \begin{quantikz}
\lstick{$x$} & \qwbundle{n} & \gate[2]{U} \gategroup[2,steps=1,style={dashed,rounded corners}]{oracle} & \rstick{$x$} \\
\lstick{$y$} & & & \rstick{$y \oplus f(x)$}
\end{quantikz} \]

Bracket tag quantikz.
\[ \begin{quantikz} & \gate{Y} & \qw \end{quantikz} \tag{Q} \]

Equation label first quantikz.
\begin{equation}
\label{eq:first}
\begin{quantikz} & \phase{\phi} & \qw \end{quantikz}
\end{equation}

\begin{theorem}
Theorem quantikz.
\begin{equation}
\begin{quantikz} & \gate{V} & \qw \end{quantikz}
\label{eq:theorem}
\end{equation}
\end{theorem}

\begin{proof}
Proof quantikz period.
\[ \begin{quantikz} & \gate{W} & \qw \end{quantikz}. \qedhere \]
\end{proof}

Figure quantikz.
\begin{figure}[h]
\centering
\begin{quantikz} & \gate{F} & \qw \end{quantikz}
\caption{A circuit in a figure}
\end{figure}

Ordinary display.
\[ |\psi\rangle = \alpha |0\rangle + \beta |1\rangle \]

Reference to the circuit: \ref{eq:circuit}, \ref{eq:first}, and \ref{eq:theorem}.
\end{document}
